% 由原 Excel 课程表转换。通常只需要编辑本文件；版式在 main.tex 中。
% 每个报告用 \ReportBlue 或 \ReportWhite，共 8 个参数：时间、报告人、题目、论文编号、辅导老师、质疑同学、提问老师、巡回点评。

\ScheduleTitle{2026-2027年秋季学期南信大方班实验班课程表}


%%%%%院士课


%方滨兴 10月29日

\Needspace{7.5cm}
\begin{tabularx}{\textwidth}{|C{1.65cm}|C{1.75cm}|Y|C{2.20cm}|C{1.70cm}|C{2.20cm}|C{1.80cm}|C{1.70cm}|}
\SessionTitle{南京信息工程大学实验班第45期课程（10月29日）报告次序表（主点评老师：待定，副点评老师：待定）\newline 腾讯会议号：991-3530-5466 线下课室：临江楼D区报告厅}
\HeaderRow

\ReportBlue{18:00-18:50}{李升冉}
{EM Eye: Characterizing Electromagnetic
Side-channel Eavesdropping on Embedded Cameras\newline EM Eye：嵌入式摄像头中电磁
侧信道窃听的特征分析}
{南信大-45-1-1}
{黄文斌}
{黄嘉乐、丁欣、袁韵涵、薛传朝、钟玉文}
{}{/}

\ReportWhite{18:50-19:40}{丁欣}
{Generative Image Steganography with  Minimum-Distance Guidance\newline 基于最小距离引导的生成式图像隐写}
{南信大-45-1-2}
{张翔}
{鲁懿、赵浩斌、罗靖昇、吴雨健、蒋承睿}
{}{/}

\BreakRow{19:40-19:50}{课间休息}

\ReportWhite{19:50-20:40}{李想}
{SELFDEFEND: LLMs Can Defend Themselves against Jailbreaking in a Practical Manner\newline }
{南信大-45-1-3}
{付章杰}
{董康、李嘉诚、周佳恩、张鑫祺、陆萍萍}
{}{方滨兴}

\ReportBlue{20:40-21:30}{赵宇阳}
{RBLJAN: Robust Byte-Label Joint Attention Network for Network Traffic Classification\newline}
{南信大-45-1-4}
{刘腾骏}
{赵德志、周蕊、朱嘉睿、仉俊杰、张耀腾}
{}{方滨兴}
\end{tabularx}
\BlockGap




%张艳宁 11月12日

\Needspace{7.5cm}
\begin{tabularx}{\textwidth}{|C{1.65cm}|C{1.75cm}|Y|C{2.20cm}|C{1.70cm}|C{2.20cm}|C{1.80cm}|C{1.70cm}|}
\SessionTitle{南京信息工程大学实验班第47期课程（11月12日）报告次序表（主点评老师：张艳宁 院士，副点评老师：待定）\newline 腾讯会议号：991-3530-5466 线下课室：临江楼D区报告厅}
\HeaderRow

\ReportBlue{18:00-18:50}{范泽楠}
{Circumventing shortcuts in audio-visual deepfake detection datasets with unsupervised learning\newline }
{南信大-47-1-1}
{宋甫元}
{周京濮、陆宇竹、孙晨明、张泽浩、谢宸亮}
{}{/}

\ReportWhite{18:50-19:40}{侯睿}
{More Than What Was Chosen: LLM-based Explainable Recommendation Beyond Noisy User Preferences\newline }
{南信大-47-1-2}
{陈先意}
{葛文武、林志鹏、郑宇博、孙孟凡、柳敏烊}
{}{/}

\BreakRow{19:40-19:50}{课间休息}

\ReportWhite{19:50-20:40}{柳敏烊}
{A Multi-Agent Collaboration via Evolving Orchestration\newline 基于演化编排的多智能体协作}
{南信大-47-1-3}
{付章杰}
{刘云飞、耿硕硕、孙祖林、蔡正、胥茗宝}
{}{/}

\ReportBlue{20:40-21:30}{陈佳杰}
{Harnessing Vision-Language Models for Time Series Anomaly Detection\newline}
{南信大-47-1-4}
{陈先意}
{王贵婵、赵泽磊、刘文海、刘孟坤、李孟豪}
{}{/}
\end{tabularx}
\BlockGap


%郭世泽 11月26日

\Needspace{7.5cm}
\begin{tabularx}{\textwidth}{|C{1.65cm}|C{1.75cm}|Y|C{2.20cm}|C{1.70cm}|C{2.20cm}|C{1.80cm}|C{1.70cm}|}
\SessionTitle{南京信息工程大学实验班第49期课程（11月26日）报告次序表（主点评老师：郭世泽 院士，副点评老师：待定）\newline 腾讯会议号：991-3530-5466 线下课室：临江楼D区报告厅}
\HeaderRow

\ReportBlue{18:00-18:50}{赵泽磊}
{SAFE RLHF: SAFE REINFORCEMENT LEARNING FROM HUMAN FEEDBACK\newline }
{南信大-49-1-1}
{姜琴}
{田嘉鑫、吴春雪、尹琨鹏、高豪博、杜佳敏}
{}{/}

\ReportWhite{18:50-19:40}{杨荣发}
{Test-Time Backdoor Detection for Object Detection Models\newline }
{南信大-49-1-2}
{闫雷鸣}
{陈谭浩嵩、王睿轩、张宇慧、高栋森、侯睿}
{}{/}

\BreakRow{19:40-19:50}{课间休息}

\ReportWhite{19:50-20:40}{罗靖昇}
{PoisonedRAG: Knowledge Corruption Attacks to Retrieval-Augmented Generation of Large Language Models\newline }
{南信大-49-1-3}
{崔琦}
{关嘉林、康梦豪、陈佳杰、赵宇阳、李慧敏}
{}{/}

\ReportBlue{20:40-21:30}{裴甲豪}
{A Wolf in Sheep’s Clothing Practical Black-box Adversarial Attacks for Evading Learning-based Windows Malware Detection in the Wild \newline 披着羊皮的狼：真实环境下针对学习型 Windows 恶意软件检测的实用黑盒逃逸攻击}
{南信大-49-1-4}
{付章杰}
{李想、孙一鸣、王少宇、石诗思、杨冬瑞}
{}{/}
\end{tabularx}
\BlockGap



%方滨兴 12月3日


\Needspace{7.5cm}
\begin{tabularx}{\textwidth}{|C{1.65cm}|C{1.75cm}|Y|C{2.20cm}|C{1.70cm}|C{2.20cm}|C{1.80cm}|C{1.70cm}|}
\SessionTitle{南京信息工程大学实验班第50期课程（12月3日）报告次序表（主点评老师：方滨兴 院士，副点评老师：待定）\newline 腾讯会议号：991-3530-5466 线下课室：临江楼D区报告厅}
\HeaderRow

\ReportBlue{18:00-18:50}{周佳恩}
{Contrastive Text-enhanced Transformer for Cross-Domain Sequential Recommendation\newline }
{南信大-50-1-1}
{王一力}
{沈涛涛、张玥宁、吕慎远、魏林东、王钇水}
{}{/}

\ReportWhite{18:50-19:40}{路怀义}
{ChatKBQA A Generate-then-Retrieve Framework for Knowledge Base Question Answering with Fine-tuned Large Language Models\newline }
{南信大-50-1-2}
{李梓强}
{刘铭濠、张继晨、张阳、周旭、徐浩洋}
{}{/}

\BreakRow{19:40-19:50}{课间休息}

\ReportWhite{19:50-20:40}{李孟豪}
{Learning Interpretable Dynamics of Stochastic Complex Systems from Experimental Data\newline }
{南信大-50-1-3}
{王一力}
{朱建华、陈丰羽、吴长凡、路怀义、陈健华}
{}{/}

\ReportBlue{20:40-21:30}{陆晗熙}
{ Catch Me if You Can: Effective Honeypot Placement in Dynamic AD Attack Graphs \newline }
{南信大-50-1-4}
{陈先意}
{范泽楠、杨荣发、杨恺杰、徐渊博、赵咨豪}
{}{/}
\end{tabularx}
\BlockGap



%王恩东：12月10日晚

\Needspace{7.5cm}
\begin{tabularx}{\textwidth}{|C{1.65cm}|C{1.75cm}|Y|C{2.20cm}|C{1.70cm}|C{2.20cm}|C{1.80cm}|C{1.70cm}|}
\SessionTitle{南京信息工程大学实验班第51期课程（12月10日）报告次序表（主点评老师：王恩东 院士，副点评老师：待定）\newline 腾讯会议号：991-3530-5466 线下课室：临江楼D区报告厅}
\HeaderRow

\ReportBlue{18:00-18:50}{徐浩洋}
{SoftCoT: Soft Chain-of-Thought for Efficient Reasoning with LLMs\newline }
{南信大-51-1-1}
{付章杰}
{李梓豪、李义辉、刘雨柯、李京、乔子龙}
{}{/}

\ReportWhite{18:50-19:40}{赵德志}
{CacheBlend: Fast Large Language Model Serving for RAG with Cached Knowledge Fusion\newline }
{南信大-51-1-2}
{李梓强}
{王运动、李升冉、陈文喆、吴世阳、颜子然}
{}{/}

\BreakRow{19:40-19:50}{课间休息}

\ReportWhite{19:50-20:40}{刘文海}
{ExpeL: LLM Agents Are Experiential Learners \newline }
{南信大-51-1-3}
{待定}
{孙焕洛、张鑫、朱毅辰、陆晗熙、黄嘉乐}
{}{/}

\ReportBlue{20:40-21:30}{过杰}
{ FedPE： Adaptive Model Pruning-Expanding forFederated Learning on Mobile Devices\newline }
{南信大-51-1-4}
{陈先意}
{丁欣、袁韵涵、薛传朝、钟玉文、鲁懿}
{}{/}
\end{tabularx}
\BlockGap



%王小云 12月17日
\Needspace{7.5cm}
\begin{tabularx}{\textwidth}{|C{1.65cm}|C{1.75cm}|Y|C{2.20cm}|C{1.70cm}|C{2.20cm}|C{1.80cm}|C{1.70cm}|}
\SessionTitle{南京信息工程大学实验班第52期（12月17日）课程报告次序表（主点评老师：王小云 院士，副点评老师：待定）\newline 腾讯会议号：991-3530-5466 线下课室：临江楼D区报告厅}
\HeaderRow

\ReportBlue{18:00-18:50}{刘雨柯}
{Is PSI Really Faster Than PSU? Achieving Efficient PSU with Invertible Bloom Filters\newline PSI 真的比 PSU 更快吗？——利用可逆布隆过滤器实现高效 PSU}
{南信大-52-1-1}
{姜琴}
{李孟豪、田嘉鑫、吴春雪、尹琨鹏、高豪博}
{}{/}

\ReportWhite{18:50-19:40}{陆宇竹}
{Audio WAtermArk: Dynamic and Harmless Watermark for Black-box Voice Dataset Copyright Protection\newline Audio WAtermArk：面向黑盒语音数据集版权保护的动态无害水印}
{南信大-52-1-2}
{何子文}
{刘铭濠、李梓豪、李义辉、刘雨柯、李京}
{}{/}

\BreakRow{19:40-19:50}{课间休息}

\ReportWhite{19:50-20:40}{郑宇博}
{AutoFHE: Automated Adaption of CNNs for Efficient Evaluation over FHE\newline AutoFHE：自动适配 CNN 以实现 FHE 上的高效评估}
{南信大-52-1-3}
{宋甫元}
{张阳、周旭、李幸、沈涛涛、张继晨}
{}{/}

\ReportBlue{20:40-21:30}{张阳}
{From Interaction to Independence: zkSNARKs for Transparent and Non-Interactive Remote Attestation\newline 从交互到独立：zkSNARKs 实现透明且非交互的远程证明}
{南信大-52-1-4}
{任勇军}
{张玥宁、吕慎远、魏林东、王钇水、杨冬瑞}
{}{/}
\end{tabularx}
\BlockGap



%云晓春：12月24日晚


\Needspace{7.5cm}
\begin{tabularx}{\textwidth}{|C{1.65cm}|C{1.75cm}|Y|C{2.20cm}|C{1.70cm}|C{2.20cm}|C{1.80cm}|C{1.70cm}|}
\SessionTitle{南京信息工程大学实验班第53期课程（12月24日）报告次序表（主点评老师：云晓春 院士，副点评老师：待定）\newline 腾讯会议号：991-3530-5466 线下课室：临江楼D区报告厅}
\HeaderRow

\ReportBlue{18:00-18:50}{周京濮}
{Backdooring Multimodal Learning \newline }
{南信大-53-1-1}
{何子文}
{赵浩斌、罗靖昇、吴雨健、蒋承睿、董康}
{}{/}

\ReportWhite{18:50-19:40}{张宇慧}
{SauronEyes Disentangling Voluminous Logs to Unveil Camouflaged Attack Intentions \newline }
{南信大-53-1-2}
{待定}
{李嘉诚、周佳恩、张鑫祺、陆萍萍、赵德志}
{}{/}

\BreakRow{19:40-19:50}{课间休息}

\ReportWhite{19:50-20:40}{刘铭濠}
{Traffic2Chain: Revealing Covert Multi-Step Attacks Through Unsupervised Traffic Behaviour Correlation \newline }
{南信大-53-1-3}
{待定}
{周蕊、朱嘉睿、仉俊杰、张耀腾、葛文武}
{}{/}

\ReportBlue{20:40-21:30}{袁韵涵}
{Weak-to-Strong Jailbreaking on Large Language Models \newline }
{南信大-53-1-4}
{闫雷鸣}
{林志鹏、郑宇博、孙孟凡、柳敏烊、周京濮}
{}{/}
\end{tabularx}
\BlockGap




%胡事民：12月31日晚


\Needspace{7.5cm}
\begin{tabularx}{\textwidth}{|C{1.65cm}|C{1.75cm}|Y|C{2.20cm}|C{1.70cm}|C{2.20cm}|C{1.80cm}|C{1.70cm}|}
\SessionTitle{南京信息工程大学实验班第54期课程（12月31日）报告次序表（主点评老师：胡事民 院士，副点评老师：待定）\newline 腾讯会议号：991-3530-5466 线下课室：临江楼D区报告厅}
\HeaderRow

\ReportBlue{18:00-18:50}{陈文喆}
{One Prompt Word is Enough to Boost Adversarial Robustness for Pre-trained Vision-Language Models \newline }
{南信大-54-1-1}
{李梓强}
{陆宇竹、孙晨明、张泽浩、谢宸亮、刘云飞}
{}{/}

\ReportWhite{18:50-19:40}{杨恺杰}
{BadThink: Triggered Overthinking Attacks on Chain-of-Thought Reasoning in Large Language Models \newline }
{南信大-54-1-2}
{熊礼治}
{耿硕硕、孙祖林、蔡正、胥茗宝、王贵婵}
{}{/}

\BreakRow{19:40-19:50}{课间休息}

\ReportWhite{19:50-20:40}{李慧敏}
{Distraction is All You Need for Multimodal Large Language Model Jailbreaking\newline }
{南信大-54-1-3}
{李梓强}
{赵泽磊、刘文海、刘孟坤、李孟豪、田嘉鑫}
{}{/}

\ReportBlue{20:40-21:30}{陈丰羽}
{Look Twice Before You Answer: Memory-Space Visual Retracing for Hallucination Mitigation in Multimodal Large Language Models \newline }
{南信大-54-1-4}
{付章杰}
{吴春雪、尹琨鹏、高豪博、杜佳敏、陈谭浩嵩}
{}{/}
\end{tabularx}
\BlockGap



%孔志印：1月7日晚




\Needspace{7.5cm}
\begin{tabularx}{\textwidth}{|C{1.65cm}|C{1.75cm}|Y|C{2.20cm}|C{1.70cm}|C{2.20cm}|C{1.80cm}|C{1.70cm}|}
\SessionTitle{南京信息工程大学实验班第55期课程（1月7日）报告次序表（主点评老师：孔志印 院士，副点评老师：待定）\newline 腾讯会议号：991-3530-5466 线下课室：临江楼D区报告厅}
\HeaderRow

\ReportBlue{18:00-18:50}{尹琨鹏}
{AdvAD: Exploring Non-Parametric Diffusion for Imperceptible Adversarial Attacks \newline }
{南信大-55-1-1}
{陈北京}
{王睿轩、张宇慧、高栋森、侯睿、关嘉林}
{}{孔志印}

\ReportWhite{18:50-19:40}{李幸}
{Red-Teaming LLM Multi-Agent Systems via Communication Attacks \newline }
{南信大-55-1-2}
{付章杰}
{张继晨、张阳、周旭、裴甲豪、沈涛涛}
{}{孔志印}

\BreakRow{19:40-19:50}{课间休息}

\ReportWhite{19:50-20:40}{李义辉}
{Class-feature Watermark: A Resilient Black-box Watermark Against Model Extraction Attacks     \newline  类别特征水印：一种抵御模型提取攻击的鲁棒黑盒水印方法}
{南信大-55-1-3}
{高光勇}
{康梦豪、陈佳杰、赵宇阳、李慧敏、李想}
{}{/}

\ReportBlue{20:40-21:30}{朱毅辰}
{VOID: Defeating Unauthorized Mimicry in Latent Diffusion Models \newline VOID：防御潜在扩散模型中的未经授权身份模仿}
{南信大-55-1-4}
{熊礼治}
{孙一鸣、王少宇、石诗思、杨冬瑞、刘铭濠}
{}{/}
\end{tabularx}
\BlockGap




%梅宏：1月14日晚



\Needspace{7.5cm}
\begin{tabularx}{\textwidth}{|C{1.65cm}|C{1.75cm}|Y|C{2.20cm}|C{1.70cm}|C{2.20cm}|C{1.80cm}|C{1.70cm}|}
\SessionTitle{南京信息工程大学实验班第56期课程（1月14日）报告次序表（主点评老师：梅宏 院士，副点评老师：待定）\newline 腾讯会议号：991-3530-5466 线下课室：临江楼D区报告厅}
\HeaderRow

\ReportBlue{18:00-18:50}{张鑫祺}
{Unlocking the Capabilities of Thought： A Reasoning Boundary Framework to Quantify and Optimize Chain-of-Thought\newline }
{南信大-56-1-1}
{待定}
{张玥宁、吕慎远、魏林东、王钇水、李梓豪}
{}{/}

\ReportWhite{18:50-19:40}{田嘉鑫}
{Chain-of-thought reasoning without prompting \newline 无需提示的思维链推理}
{南信大-56-1-2}
{待定}
{李义辉、刘雨柯、李京、乔子龙、朱建华}
{}{/}

\BreakRow{19:40-19:50}{课间休息}

\ReportWhite{19:50-20:40}{仉俊杰}
{REGENT: A Retrieval-Augmented Generalist Agent That Can Act In-Context in New Environments \newline }
{南信大-56-1-3}
{崔琦}
{陈丰羽、吴长凡、路怀义、陈健华、范泽楠}
{}{/}

\ReportBlue{20:40-21:30}{林志鹏}
{MemAgent: Reshaping Long-Context LLM with Multi-Conv RL-based Memory Agent \newline }
{南信大-56-1-4}
{王晓苏}
{杨荣发、杨恺杰、徐渊博、赵咨豪、王运动}
{}{/}
\end{tabularx}
\BlockGap



